\documentclass[a4paper, twoside, titlepage, 12pt]{article}
% Comenzamos a añadir paquetes:

\usepackage{polyglossia}
\setmainlanguage{spanish}
%\setdefaultlanguage{galician}
% \setdefaultlanguage{galician} % Esto es idéntico a escribir la linea anterior ¯\_(ツ)_/¯
% Cambia cosas típicas de cada idioma como: «chapter», «section», «figure», «author»... por «capítulo», «sección», «figura», «autor»...

% hey

\usepackage{fontspec}
%\setmainfont[Ligatures=TeX]{Linux Biolinum O}
% \setmainfont[Ligatures=TeX]{Linux Libertine O} % El nombre de la fuente en OS X.
\setmainfont[Ligatures=TeX]{Linux Libertine} % El nombre de la fuente en Windows.

\usepackage{setspace}
\spacing{1.5}

\usepackage{blindtext} % Para generar texto de relleno.

\usepackage{tocloft} % Para tunear la tabla de contenidos.

% Este paquete incluye los logos de XeLaTeX y otros.
\usepackage{metalogo}

% Este paquete vale para hacer tablas estilo libro "serio".
\usepackage{booktabs}

% Este paquete es el que utilizamos para incluir imágenes en el documento.
\usepackage{graphicx}

% Este paquete se emplea para la gestión de hyperlinks dentro del documento: el índice apunta a cada capítulo, sección, etc., insertar urls externas dentro del documento, etc.
%\usepackage[hyperref, xetex, svgnames]{xcolor} % Distintos nombres para colores.
\usepackage[hyperref, xetex, dvipsnames]{xcolor} % svgnames ofrence uno cto. de nombres y dvipsnames otro.
\usepackage{hyperref}
\hypersetup{ % Nos permite especificar los colores con que se mostrarán los distintos tipos de link dentro del documento.
    colorlinks=true,
    linkcolor=black,
    filecolor=magenta,
    urlcolor=cyan,
    citecolor=BlueViolet
}

%\usepackage{biblatex} % Bibliografía
%\addbibresource{biblio.bib} % Archivo con las entradas de bibliografía.
\usepackage{apacite} % Bibliografía con normas APA.

% Este paquete permite que las tablas tengan gran cantidad de texto en una celda (entre otras muchas cosas)
\usepackage{tabularx}

\usepackage{titling} % Para hacer referencia al título y al autor en la portada.

% Esto se usará en la página del título (al usar \maketitle).
\title{Mi mini manual de \XeLaTeX}
\author{Manuel C. Piñeiro Mourazos}

\begin{document}
\begin{titlepage}
  \centering
	\includegraphics[width=0.25\textwidth]{images/escudo_uvigo}\par\vspace{1cm}
	\includegraphics[width=\textwidth]{images/Logotipo_negro_UNIVERSIDADE_DE_VIGO}\par\vspace{1cm}
	{\scshape\Large Pre-trabajo de Fin de Grado\par}
	\vspace{1.5cm}
	{\huge\bfseries Educación con Evidencias\par}
	\vspace{2cm}
	{\Large\itshape \theauthor\par}
	\vfill
% Bottom of the page
	{\large \today\par}
\end{titlepage}

\tableofcontents % La tabla de contenidos
%\listoffigures % La lista de imágenes (descomentar si tienes y quieres que salgan).
%\listoftables % La lista de tablas (descomentar si tienes y quieres que salgan).
\newpage %Después de las tablas de contenidos, figuras, etc. empezar en una nueva página.

\section{Introducción}
Título del trabajo de fin de grado.

\section{Situación actual}
Situación actual más antecedentes.

\section{Analizar las causas}
Situación acutal más antecedentes.

\section{Objetivo}
Objetivo general y objetivos específicos.

\section{Plan de acción}
Metodología según el tipo de trabajo elegido.

\section{Seguimiento}
Resultado, evaluación.

\section{Resultados obtenidos}
Conclusiones.

\end{document}
